\begin{table}
\centering
\caption{\label{tab:qp_cold_wave_city}cold wave model}
\centering
\resizebox{\ifdim\width>\linewidth\linewidth\else\width\fi}{!}{
\begin{tabu} to \linewidth {>{\raggedright}X>{\raggedright}X>{\raggedright}X>{\raggedright}X>{\raggedright}X}
\toprule
  & minimum & mean & maximum & dew point\\
\midrule
cold wave & 0.322 & 0.106 & 0.267 & 0.013\\

 & (0.208) & (0.262) & (0.258) & (0.429)\\

Num.Obs. & 125343 & 125343 & 125344 & 125368\\
\bottomrule
\multicolumn{5}{l}{\rule{0pt}{1em}\textit{Note: }}\\
\multicolumn{5}{l}{\rule{0pt}{1em}+ p < 0.1, * p < 0.05, ** p < 0.01, *** p < 0.001}\\
\multicolumn{5}{l}{\rule{0pt}{1em}All models include city-by-year, day-of-week}\\
\multicolumn{5}{l}{\rule{0pt}{1em}and month-of-year fixed effects.}\\
\multicolumn{5}{l}{\rule{0pt}{1em}Standard errors are clustered at the city level.}\\
\multicolumn{5}{l}{\rule{0pt}{1em}Temperature thresholds reflect the city-level}\\
\multicolumn{5}{l}{\rule{0pt}{1em}1textsuperscript{st} percentile of the distribution.}\\
\end{tabu}}
\end{table}
